\section{Aplicaciones y limitaciones de las RNN}

\subsection{Un caso de aplicación de las RNN: generación de música}

El presentador presentó la aplicación de las RNN en la generación de música. Esta es una tarea naturalmente adecuada para el modelado de secuencias: dada la secuencia histórica de notas musicales, se predice la nota siguiente más probable.

\begin{figure}[H]
\centering
\includegraphics[width=0.85\textwidth]{slides-images/slide-56.jpg}
\caption{Generación de música con RNN: predicción de la siguiente nota\protect\footnotemark}
\end{figure}
\footnotetext{Diapositivas, página 56.}

El presentador mencionó un caso interesante: un equipo entrenó un modelo de red neuronal para aprender música clásica y después le pidió al modelo que intentara completar el tercer movimiento de la célebre «Sinfonía inconclusa» de Franz Schubert. Esto muestra el potencial de las RNN en tareas creativas.

\subsection{Criterios de diseño del modelado de secuencias}

El presentador resumió cuatro criterios de diseño clave del modelado de secuencias:

\begin{figure}[H]
\centering
\includegraphics[width=0.6\textwidth]{slides-images/slide-30.jpg}
\caption{Criterios de diseño del modelado de secuencias\protect\footnotemark}
\end{figure}
\footnotetext{Diapositivas, página 30.}

\begin{enumerate}
    \item \textbf{Manejar secuencias de longitud variable} (Variable-length sequences)
    \item \textbf{Rastrear dependencias de largo plazo} (Long-term dependencies)
    \item \textbf{Mantener información sobre el orden} (Maintain information about order)
    \item \textbf{Compartir parámetros a lo largo de la secuencia} (Share parameters across the sequence)
\end{enumerate}

Las RNN satisfacen estos cuatro criterios, pero en la práctica siguen teniendo limitaciones importantes.

\subsection{Limitaciones centrales de las RNN}

\begin{warningbox}{Las tres grandes limitaciones de las RNN}
\begin{enumerate}
    \item \textbf{Cuello de botella del cómputo secuencial}: las RNN deben procesar los datos paso a paso en el tiempo y no pueden paralelizarse. Para secuencias largas, tanto el entrenamiento como la inferencia son muy lentos
    \item \textbf{Dificultad con las dependencias de largo plazo}: aunque las LSTM mitigan el problema del desvanecimiento del gradiente, para secuencias muy largas la información sigue atenuándose de manera gradual
    \item \textbf{Cuello de botella de codificación}: en la arquitectura Encoder-Decoder, la información de toda la secuencia de entrada debe comprimirse en un vector de estado oculto de tamaño fijo, lo cual resulta insuficiente para secuencias largas
\end{enumerate}
\end{warningbox}

Estas limitaciones llevaron a los investigadores a preguntarse: \textbf{¿es posible prescindir por completo del mecanismo de recurrencia (Recurrence) y procesar los datos de secuencia de una manera totalmente nueva?}

\subsection{Resumen de la sección}

Las RNN (incluidas las LSTM/GRU) lograron un gran éxito en el modelado de secuencias, pero su naturaleza de procesamiento secuencial limita la eficiencia computacional, y las dificultades en la propagación del gradiente también restringen su capacidad para manejar secuencias extremadamente largas. Estas limitaciones impulsaron directamente el surgimiento del mecanismo de Attention y de la arquitectura Transformer.

\newpage
